%============================================================
% CH2 · PATTERN 2 — EXERCISES, HINTS, ANSWERS
%============================================================
\subsection{Pattern 2: Exercises and Going Further}
\label{ssec:ch2-p2-ex}

\begin{patternbox}{Pattern 2 Summary Card}
\renewcommand{\arraystretch}{1.25}
\begin{tabularx}{\linewidth}{@{}>{\bfseries}l X@{}}
Recognise & ``for every vertex, \dots{} in its subtree'' (up); ``\dots{} along the path from the root'', ``add to a whole subtree'' (down); ``delete a vertex / cut an edge'' (pieces) \\
Root it & pass the parent instead of \texttt{visited}; never trust the orientation of the input edges \\
Up & at leave: $\mathrm{agg}[v] = \mathrm{base}(v) \oplus \mathrm{agg}[\text{children}]$; the part above has $n - \mathrm{sz}[v]$ vertices \\
Down & at enter: $\mathrm{acc}[v] = \mathrm{acc}[\mathrm{par}] \otimes \mathrm{step}(v)$; for subtree updates mark the top and push down once \\
No recursion & BFS order: forward $=$ down, backward $=$ up; parent array with $p_i < i$: the labels are the order \\
Top traps & wrong parent from $a_i < b_i$; a value used before it is final; forgetting $n - \mathrm{sz}$; \texttt{int} products; deep recursion \\
\end{tabularx}
\end{patternbox}

\exlegend
\begin{enumerate}[leftmargin=*, label=\textbf{E2.\arabic*}, itemsep=6pt]
    \item \exeasy\ \pconcept \lc{Medium}\ \exprob{Number of Nodes in the Sub-Tree With the Same Label}{LeetCode 1519}{https://leetcode.com/problems/number-of-nodes-in-the-sub-tree-with-the-same-label/}. A tree on $0..n-1$ rooted at 0, each node labelled with a lowercase letter. For every node $i$, return how many nodes of its subtree carry the same label as $i$. ($n \le 10^5$.)
    \item \exeasy\ \pclassic \cf{1500}\ \exprob{Cut 'em all!}{Codeforces 982C}{https://codeforces.com/problemset/problem/982/C}. Remove as many edges of a tree as possible so that every remaining component has an even number of vertices; print that number, or $-1$ if impossible. ($n \le 10^5$.)
    \item \exmedium\ \pclassic \lc{Medium}\ \exprob{Minimum Fuel Cost to Report to the Capital}{LeetCode 2477}{https://leetcode.com/problems/minimum-fuel-cost-to-report-to-the-capital/}. A tree of cities with capital 0; one representative and one car with \texttt{seats} seats in every city; driving a road costs one litre per car; representatives may change cars. Minimum total fuel for everybody to reach the capital. ($n, \texttt{seats} \le 10^5$.)
    \item \exmedium\ \ptricky \cf{1600}\ \exprob{Alyona and the Tree}{Codeforces 682C}{https://codeforces.com/problemset/problem/682/C}. A tree rooted at 1 with a number $a_u$ on every vertex and a (possibly negative) number on every edge. Vertex $v$ is \emph{sad} if some $u$ in its subtree has $\mathrm{dist}(v, u) > a_u$. Remove leaves one at a time until no vertex is sad; minimise the number removed. ($n \le 10^5$, $|c| \le 10^9$.) The trap: the edge weights can be negative, so the worst ancestor is not the root.
    \item \exmedium\ \pconcept \lc{Hard}\ \exprob{Maximum Number of K-Divisible Components}{LeetCode 2872}{https://leetcode.com/problems/maximum-number-of-k-divisible-components/}. Remove any set of edges of a tree with node values so that every component's value sum is divisible by $k$ (the total is). Maximise the number of components. ($n \le 3\cdot10^4$.)
    \item \exhard\ \precent \lc{Hard}\ \exprob{Find Number of Coins to Place in Tree Nodes}{LeetCode 2973}{https://leetcode.com/problems/find-number-of-coins-to-place-in-tree-nodes/}. A tree rooted at 0 with a non-zero cost on every node. For each node: 1 coin if its subtree has fewer than 3 nodes, otherwise the maximum product of the costs of 3 distinct nodes of its subtree, or 0 if that product is negative. ($n \le 2\cdot10^4$, $|\mathrm{cost}| \le 10^4$.)
\end{enumerate}

\subsubsection*{Hints}
\begin{enumerate}[leftmargin=*, label=\textbf{E2.\arabic*}]
    \item Merging 26 counters per node works, but there is a cheaper way: a \emph{single} global counter per letter. What does it hold when you enter $i$, and when you leave $i$?
    \item With $n$ odd the total is odd: $-1$. Otherwise, which edges can be cut on their own? Does cutting one change whether another can be cut?
    \item Look at one road: how many representatives must cross it, at least? Can all roads reach their lower bound at the same time?
    \item A vertex $u$ that makes some ancestor sad must go, and with it its whole subtree. Which ancestor maximises $\mathrm{dist}(v, u)$? Think of the best upward path ending at $u$.
    \item Process children before parents. When the remaining sum of a subtree is divisible by $k$, is it ever wrong to cut it off?
    \item A maximum product of three numbers uses either the three largest, or the two smallest (both negative) and the largest. How many numbers must a subtree pass to its parent?
\end{enumerate}

\subsubsection*{Answers}

\begin{solutionbox}[E2.1 Number of Nodes in the Sub-Tree With the Same Label]
\textbf{The before/after trick.} Keep one global counter per letter, incremented when a node is \emph{entered}. Between entering and leaving $i$, exactly the nodes of $i$'s subtree are entered (Lesson~\ref{ssec:ch2-call-stack}), so $\mathrm{ans}[i] = \mathrm{cnt}_{\text{leave}}[\ell_i] - \mathrm{cnt}_{\text{enter}}[\ell_i]$, where the entry value is read \emph{before} $i$ itself is counted. This needs a real DFS (the BFS order does not keep subtrees together), so it is the iterative DFS with \texttt{it[u]}; $n = 10^5$ rules out recursion on LeetCode.
\begin{center}
\begin{tikzpicture}[x=1.15cm, y=0.9cm]
    \node[gvsrc] (a0) at (0,0) {0}; \node[gv] (a1) at (-1,-1) {1}; \node[gv] (a3) at (1,-1) {3}; \node[gv] (a2) at (-1,-2) {2};
    \foreach \a/\b in {0/1, 1/2, 0/3} \draw[gtree] (a\a) -- (a\b);
    \node[note, anchor=west, text width=8.2cm, align=left] at (2.2,-1) {labels \texttt{bbbb}. Counter of \texttt{b}: enter 0 (before 0), enter 1 (before 1), enter 2 (before 2), leave 2 at 3 $\Rightarrow 1$, leave 1 at 3 $\Rightarrow 2$, enter 3 (before 3), leave 3 at 4 $\Rightarrow 1$, leave 0 at 4 $\Rightarrow 4$};
\end{tikzpicture}
\end{center}
\begin{lstlisting}[style=cpp]
#include <bits/stdc++.h>
using namespace std;

class Solution {
public:
    vector<int> countSubTrees(int n, vector<vector<int>>& edges, string labels) {
        vector<vector<int>> adj(n);
        for (auto& e : edges) { adj[e[0]].push_back(e[1]); adj[e[1]].push_back(e[0]); }
        vector<int> ans(n, 0), before(n, 0), parent(n, -1), cnt(26, 0);
        vector<size_t> it(n, 0);
        vector<int> st(1, 0);
        before[0] = cnt[labels[0] - 'a']++;            // enter 0: remember the old count
        while (!st.empty()) {
            int u = st.back();
            if (it[u] < adj[u].size()) {
                int v = adj[u][it[u]++];
                if (v == parent[u]) continue;
                parent[v] = u;
                before[v] = cnt[labels[v] - 'a']++;    // enter v
                st.push_back(v);
            } else {                                   // leave u: its whole subtree was counted
                ans[u] = cnt[labels[u] - 'a'] - before[u];
                st.pop_back();
            }
        }
        return ans;
    }
};
\end{lstlisting}
Examples: $[2,1,1,1,1,1,1]$, $[4,2,1,1]$, $[3,2,1,1,1]$. On a path of $10^5$ nodes all labelled \texttt{z} it returns $10^5$ for the root and 1 for the last node. Matched a brute force (walk every subtree and count) on $5000$ random trees with $n \le 12$ and 3 letters. The alternative --- an array of 26 counts per node, added into the parent in the backward BFS loop --- is also correct and \bigO{26n}, but needs $26n$ integers of memory; the trick above needs 26 in total.
\end{solutionbox}

\begin{solutionbox}[E2.2 Cut 'em all!]
If $n$ is odd, some component is odd: $-1$. Otherwise cut the edge above $v$ exactly when $\mathrm{sz}[v]$ is even. \emph{Why this is allowed:} cutting such an edge leaves an even piece below and an even remainder above, so parities elsewhere do not change --- every such edge can be cut independently. \emph{Why no other edge:} after any valid set of cuts, the vertices below the edge above $v$ form a union of whole components iff that edge is cut; a union of even components has even size, so a cut edge always has $\mathrm{sz}[v]$ even. Hence the answer is the number of non-root $v$ with even $\mathrm{sz}[v]$.
\begin{center}
\begin{tikzpicture}[x=1.2cm, y=0.9cm]
    \node[gvsrc] (c1) at (0,0) {1}; \node[gv] (c4) at (-0.8,-1) {4}; \node[gv] (c3) at (0.8,-1) {3}; \node[gv] (c2) at (-0.8,-2) {2};
    \draw[gtree] (c1) -- (c3); \draw[gtree] (c4) -- (c2); \draw[gcut] (c1) -- (c4);
    \node[gdist, text=primaryblue, left=1pt of c4] {sz 2}; \node[gdist, text=primaryblue, right=1pt of c3] {sz 1}; \node[gdist, text=primaryblue, left=1pt of c2] {sz 1};
    \node[note, anchor=west, text width=6cm, align=left] at (2.0,-1) {sample 1, rooted at 1: only $\mathrm{sz}[4] = 2$ is even $\Rightarrow$ cut $1{-}4$, answer 1};
\end{tikzpicture}
\end{center}
\begin{lstlisting}[style=cpp]
#include <bits/stdc++.h>
using namespace std;

// Cut 'em all! (CF 982C): an edge (parent, v) can be cut iff sz[v] is even; with n odd nothing works.
int main() {
    int n;
    cin >> n;
    vector<vector<int>> adj(n + 1);
    for (int i = 0; i < n - 1; i++) {
        int a, b;
        cin >> a >> b;
        adj[a].push_back(b);
        adj[b].push_back(a);
    }
    if (n % 2 == 1) { cout << -1 << "\n"; return 0; }
    vector<int> parent(n + 1, 0), order(1, 1), sz(n + 1, 1);
    for (size_t i = 0; i < order.size(); i++)
        for (int v : adj[order[i]])
            if (v != parent[order[i]]) { parent[v] = order[i]; order.push_back(v); }
    int cuts = 0;
    for (int i = n - 1; i >= 1; i--) {             // reverse BFS order: children first
        int v = order[i];
        if (sz[v] % 2 == 0) cuts++;                // sz[v] is final: cut the edge above v
        sz[parent[v]] += sz[v];
    }
    cout << cuts << "\n";
}
\end{lstlisting}
Samples: \texttt{1}, \texttt{-1}, \texttt{4}, \texttt{0}. Matched a brute force that tries all $2^{n-1}$ edge subsets on $3000$ random trees with $n \le 12$ (including $n = 1$, answer $-1$).
\end{solutionbox}

\begin{solutionbox}[E2.3 Minimum Fuel Cost to Report to the Capital]
\textbf{Lower bound.} The road above $v$ must be crossed by all $\mathrm{sz}[v]$ representatives of the subtree of $v$, so it is driven by at least $\lceil \mathrm{sz}[v] / \texttt{seats} \rceil$ cars. \textbf{Achievable.} Move everybody up level by level: at $v$, the $\mathrm{sz}[v]$ people present pack into $\lceil \mathrm{sz}[v]/\texttt{seats} \rceil$ cars --- there are enough cars, because every person arrived with a car or in one ($\ge$ that many cars are at $v$). So the answer is $\sum_{v \ne 0} \lceil \mathrm{sz}[v] / \texttt{seats} \rceil$: aggregation plus one term per edge.
\begin{lstlisting}[style=cpp]
#include <bits/stdc++.h>
using namespace std;
typedef long long ll;

class Solution {
public:
    long long minimumFuelCost(vector<vector<int>>& roads, int seats) {
        int n = roads.size() + 1;
        vector<vector<int>> adj(n);
        for (auto& r : roads) { adj[r[0]].push_back(r[1]); adj[r[1]].push_back(r[0]); }
        vector<int> parent(n, -1), order(1, 0);
        for (size_t i = 0; i < order.size(); i++)
            for (int v : adj[order[i]])
                if (v != parent[order[i]] && v != 0) { parent[v] = order[i]; order.push_back(v); }
        vector<ll> sz(n, 1);
        ll fuel = 0;
        for (int i = n - 1; i >= 1; i--) {
            int v = order[i];
            fuel += (sz[v] + seats - 1) / seats;       // sz[v] people cross the road above v
            sz[parent[v]] += sz[v];
        }
        return fuel;
    }
};
\end{lstlisting}
Examples: \textbf{3}, \textbf{7}, \textbf{0}. The answer needs 64 bits: a path of $10^5$ cities with one seat costs $\sum_{k=1}^{99\,999} k = 4\,999\,950\,000$ (printed by the program). Checked on $5000$ random trees ($n \le 14$, $\texttt{seats} \le 4$) against an independent computation of the crossings (every representative walks its own path to the capital) and against a simulation that moves the cars level by level and asserts that enough cars are present. The simulation verifies feasibility; optimality is the lower bound above.
\end{solutionbox}

\begin{solutionbox}[E2.4 Alyona and the Tree]
\textbf{Which vertices must go.} If $u$ makes an ancestor $v$ sad, $u$ cannot stay (removing leaves never shortens the path $v \leadsto u$), and removing $u$ means removing its whole subtree first. Conversely, after removing exactly the subtrees of such ``bad'' $u$, nobody is sad. So the answer is the number of vertices with a bad vertex on their root path (itself included).

\textbf{The trap.} $\max_v \mathrm{dist}(v, u)$ over ancestors $v$ is \emph{not} $\mathrm{dist}(1, u)$: with negative edges, a later start can be better. It is the best sum of an upward path ending at $u$, a push-down with reset (Kadane's algorithm along every root path): $\mathrm{best}[u] = \max(\mathrm{best}[\mathrm{par}], 0) + c(\mathrm{par}, u)$. The ``$\max(\cdot, 0)$'' starts a fresh path at the parent. Do not descend below a bad vertex; the answer is $n - $ (vertices reached). Values reach $10^5 \cdot 10^9$: \texttt{long long}.
\begin{lstlisting}[style=cpp]
#include <bits/stdc++.h>
using namespace std;
typedef long long ll;

// Alyona and the Tree (CF 682C). best[u] = max over ancestors v of dist(v, u) = the largest sum of a
// path that ENDS at u and goes up: best[u] = max(best[p], 0) + c(p, u). A vertex with best[u] > a[u]
// must go, and with it its whole subtree. Answer = n - (vertices kept).
int main() {
    ios::sync_with_stdio(false);
    cin.tie(nullptr);
    int n;
    cin >> n;
    vector<ll> a(n + 1);
    for (int i = 1; i <= n; i++) cin >> a[i];
    vector<vector<pair<int, ll>>> adj(n + 1);
    for (int i = 2; i <= n; i++) {
        int p; ll c;
        cin >> p >> c;
        adj[p].push_back({i, c});
        adj[i].push_back({p, c});
    }
    vector<int> parent(n + 1, 0), order(1, 1);
    vector<ll> best(n + 1, 0);
    int kept = 0;
    for (size_t i = 0; i < order.size(); i++) {    // BFS = pre-order push-down, no recursion
        int u = order[i];
        kept++;
        for (auto& e : adj[u]) {
            int v = e.first;
            if (v == parent[u]) continue;
            parent[v] = u;
            best[v] = max(best[u], 0LL) + e.second; // restart the path at u if that is better
            if (best[v] > a[v]) continue;           // v makes an ancestor sad: drop v and its subtree
            order.push_back(v);
        }
    }
    cout << n - kept << "\n";
}
\end{lstlisting}
Sample: \textbf{5}. Matched a brute force (for every $u$, walk all ancestors and test $\mathrm{dist}(v,u) > a_u$; then count vertices below a bad one) on $3000$ random trees with $n \le 10$, edge weights in $[-20, 20]$ and parents that may carry larger labels than their children --- the input lists ``edge between $i+1$ and $p_i$'', and $p_i$ is the parent but not necessarily smaller.
\end{solutionbox}

\begin{solutionbox}[E2.5 Maximum Number of K-Divisible Components]
Children first (backward BFS order), keep $\mathrm{rem}[v]$ = the sum of what is still attached to $v$, modulo $k$. If $\mathrm{rem}[v] = 0$ when $v$ is left, cut the edge above $v$ (for the root: the last component); otherwise hand $\mathrm{rem}[v]$ to the parent. \emph{Why greedy is safe:} if the attached part of the subtree of $v$ has sum $\equiv 0$, cutting it off changes neither the residue of anything above (it removes $0 \bmod k$) nor any choice below; and an edge above a $v$ with $\mathrm{rem}[v] \ne 0$ can never be cut, because the part below it would have a sum $\not\equiv 0$. So the greedy cuts every edge that any valid split could cut.
\begin{center}
\begin{tikzpicture}[x=1.2cm, y=0.9cm]
    \node[gvsrc] (k0) at (0,0) {0}; \node[gv] (k2) at (0,-1) {2}; \node[gv] (k1) at (-0.9,-2) {1}; \node[gv] (k4) at (0.9,-2) {4}; \node[gv] (k3) at (-0.9,-3) {3};
    \draw[gtree] (k0) -- (k2); \draw[gtree] (k2) -- (k4); \draw[gtree] (k1) -- (k3); \draw[gcut] (k2) -- (k1);
    \foreach \v/\x in {0/1, 2/1, 1/8, 4/4, 3/4} \node[gdist, text=deepblue, right=1pt of k\v] {\x};
    \node[note, anchor=west, text width=6.4cm, align=left] at (2.0,-1.5) {example 1, $k = 6$ (values in blue). Subtree of 1: $8 + 4 = 12 \equiv 0$: cut. Node 2: $1 + 4 = 5$, handed up; root: $1 + 5 = 6 \equiv 0$: the last piece. Answer 2};
\end{tikzpicture}
\end{center}
\begin{lstlisting}[style=cpp]
#include <bits/stdc++.h>
using namespace std;
typedef long long ll;

class Solution {
public:
    int maxKDivisibleComponents(int n, vector<vector<int>>& edges, vector<int>& values, int k) {
        vector<vector<int>> adj(n);
        for (auto& e : edges) { adj[e[0]].push_back(e[1]); adj[e[1]].push_back(e[0]); }
        vector<int> parent(n, -1), order(1, 0);
        for (size_t i = 0; i < order.size(); i++)
            for (int v : adj[order[i]])
                if (v != parent[order[i]] && v != 0) { parent[v] = order[i]; order.push_back(v); }
        vector<ll> rem(n);                             // subtree sum mod k (of what is still attached)
        for (int v = 0; v < n; v++) rem[v] = values[v] % k;
        int comps = 0;
        for (int i = n - 1; i >= 0; i--) {
            int v = order[i];
            if (rem[v] == 0) comps++;                  // cut above v (the root: the last piece)
            else rem[parent[v]] = (rem[parent[v]] + rem[v]) % k;
        }
        return comps;
    }
};
\end{lstlisting}
Examples: \textbf{2}, \textbf{3}. Matched a brute force over all $2^{n-1}$ edge subsets on $5000$ random trees ($n \le 11$, $k \le 5$, values $0..6$ adjusted so that the total is divisible by $k$). Reducing modulo $k$ at every step keeps the numbers small; without it the sums still fit in \texttt{long long} ($3\cdot10^4 \cdot 10^9$), but not in \texttt{int}.
\end{solutionbox}

\begin{solutionbox}[E2.6 Find Number of Coins to Place in Tree Nodes]
The best product of three numbers from a multiset is $\max(\text{top}_1\text{top}_2\text{top}_3,\ \text{low}_1\text{low}_2\,\text{top}_1)$: either the three largest, or the two most negative times the largest. So a subtree needs to pass up only its \textbf{3 largest and 2 smallest} costs --- at most 5 numbers. Each node merges its own cost with its children's lists, sorts, drops everything between the 2nd smallest and the 3rd largest, and evaluates the formula if its subtree has $\ge 3$ nodes. Since every child passes at most 5 numbers, node $v$ sorts at most $5\deg(v) + 1$ numbers: \bigO{n \log n} in the worst case (a star), \bigO{n} on bounded degrees.
\begin{lstlisting}[style=cpp]
#include <bits/stdc++.h>
using namespace std;
typedef long long ll;

class Solution {
public:
    vector<long long> placedCoins(vector<vector<int>>& edges, vector<int>& cost) {
        int n = cost.size();
        vector<vector<int>> adj(n);
        for (auto& e : edges) { adj[e[0]].push_back(e[1]); adj[e[1]].push_back(e[0]); }
        vector<int> parent(n, -1), order(1, 0);
        for (size_t i = 0; i < order.size(); i++)
            for (int v : adj[order[i]])
                if (v != parent[order[i]] && v != 0) { parent[v] = order[i]; order.push_back(v); }
        // keep[v] = at most 5 costs of v's subtree: its 3 largest and its 2 smallest
        vector<vector<int>> keep(n);
        vector<int> sz(n, 1);
        vector<ll> coin(n, 1);
        for (int v = 0; v < n; v++) keep[v].push_back(cost[v]);
        for (int i = n - 1; i >= 0; i--) {
            int v = order[i];
            vector<int>& c = keep[v];                  // all children of v were merged into c
            sort(c.begin(), c.end());
            if (c.size() > 5) c.erase(c.begin() + 2, c.end() - 3);   // drop the middle
            if (sz[v] >= 3) {
                int m = c.size();
                ll top3 = (ll)c[m - 1] * c[m - 2] * c[m - 3];
                ll two_neg = (ll)c[0] * c[1] * c[m - 1];
                coin[v] = max(0LL, max(top3, two_neg));
            }
            if (i > 0) {
                int p = parent[v];
                sz[p] += sz[v];
                keep[p].insert(keep[p].end(), c.begin(), c.end());
            }
        }
        return coin;
    }
};
\end{lstlisting}
Examples: $[120,1,1,1,1,1]$, $[280,140,32,1,1,1,1,1,1]$, $[0,1,1]$. Matched a brute force that collects every subtree and tries all triples on $5000$ random trees ($n \le 13$, costs in $\pm[1, 9]$). Two details: with only 3 or 4 kept values the ``top'' and ``low'' entries overlap, which is fine --- both formulas still use three different positions of the sorted list, i.e.\ three distinct nodes; and the products reach $10^{12}$, so they are computed in \texttt{long long}.
\end{solutionbox}
